\section{Background and Related Work}
\label{sec:related}

\subsection{Dealer Hedging and Market Microstructure}

Market makers face regulatory requirements to maintain delta-neutral positions, creating systematic hedging flows when gamma exposure accumulates. \citet{grossman1988liquidity} demonstrated how market maker inventory risk creates predictable hedging patterns, while \citet{frey1997market} formalized how dynamic hedging amplifies volatility through feedback effects.

A critical principle underlying our methodology is the dealer counterparty framework: market makers serve as counterparties to customer option positions, such that dealer gamma equals the negative of aggregate customer gamma. \citet{anderegg2022impact} establish this relationship empirically using trade-repository data, demonstrating that aggregated market makers carry negative gamma exposure that feeds back to the spot market; \citet{adams2025zerodte} extend the framework to the 0DTE setting by reconstructing dealer hedging needs from intraday order flow.

Empirically, \citet{ni2005stock} documented stock price clustering on option expiration dates, providing evidence that dealer hedging creates measurable price effects. \citet{garleanu2009demand} developed a demand-based option pricing model showing how dealer constraints affect both option prices and underlying stock dynamics. Practitioner research has operationalized gamma exposure through standardized GEX calculations \citep{spotgamma2021,anderegg2022impact}, though these aggregate scalar metrics function as lossy compression of the underlying strike-level distribution---a limitation our raw chain validation directly addresses (Section~\ref{subsec:raw_chain}).

\subsection{Zero-Days-to-Expiration (0DTE) Options}

The introduction of daily options expirations in 2022 fundamentally altered market structure \citep{cboe2024zero}. By 2023, 0DTE options accounted for 43\% of SPX options volume (up from $<$5\% in 2020), concentrating gamma exposure at very short maturities. \citet{dim2023odtes} provide the first systematic empirical study of 0DTE dealer inventory: market-maker net gamma is on average positive and negatively related to future intraday volatility, with delta-hedging-consistent price dynamics that are inconsistent with information-based trading, establishing dealer-hedging rather than information flow as the dominant channel through which 0DTE trading affects the underlying. \citet{adams2025zerodte} extend this empirical characterisation by reconstructing market-maker hedging needs from intraday flow and documenting that 0DTE proliferation has fundamentally restructured dealer hedging dynamics, with positions in longer-dated options that subsequently become 0DTE driving the bulk of the structural shift. Our 2022--2025 multi-year panel (Section~\ref{sec:regime}) is consistent with this characterisation: detection of persistent dealer-gamma regimes grows from 3.7\% of 30-day windows in 2021 to 100\% in 2024--2025, coincident with but not a proof of the 0DTE mechanism.

This structural shift has two consequences relevant to our framework. First, daily 0DTE rollovers may create more sustained directional gamma exposure than traditional monthly expiration cycles, as each day's new contracts concentrate gamma at near-the-money strikes before decaying to zero. Second, the sheer volume concentration means dealer hedging flows from 0DTE contracts dwarf those from traditional options, potentially creating persistent negative gamma environments where regime-switching dynamics---historically the norm \citep{ni2005stock}---give way to sustained directional positioning.

\subsection{LLM Reasoning Evaluation}

Evaluating LLM reasoning capabilities beyond surface-level pattern matching is an active research area. \citet{mccoy2023embers} demonstrate that LLMs are fundamentally shaped by training probability distributions, performing better on high-probability sequences---raising concerns about memorization versus genuine reasoning. Chain-of-thought prompting \citep{openai2024reasoning} enables transparent reasoning traces but does not inherently prevent training data leakage. \citet{wei2022chain} demonstrated that chain-of-thought enables complex multi-step reasoning, while \citet{kojima2022large} showed zero-shot reasoning without task-specific training. However, \citet{marcus2019rebooting} argues that distinguishing true reasoning from memorization remains challenging.

\subsection{LLM Applications in Finance}

LLMs have shown promise in financial analysis: \citet{lopez2023can} demonstrate that LLMs can extract market-relevant information from financial text. However, \citet{lopezlira2025memorization} document a parallel risk---that LLMs may rely on memorized economic facts from their training period rather than genuine reasoning, with this memorization persisting despite explicit instructions to respect historical boundaries. More broadly, \citet{dong2024generalization} establish that data contamination corrupts LLM benchmark performance and propose detection methodologies, but no comparable framework exists for validating structural reasoning in market microstructure. Recent work also employs LLMs as trading agents \citep{yang2024tradingagents}, though validation of structural reasoning---rather than pattern recall---remains limited. Two critical challenges therefore persist: \textit{temporal memorization} (recalling specific events from training data) and \textit{spurious pattern detection} (identifying statistically significant but mechanically meaningless patterns). Our obfuscation framework addresses both.

\subsection{Regime Detection in Financial Markets}

Regime detection has a rich history in financial econometrics. \citet{hamilton1989new} introduced Markov-switching models for identifying volatility regimes from return series. \citet{ang2002regime} applied regime-switching to asset pricing, demonstrating that accounting for regime shifts improves portfolio allocation and risk management. More recently, \citet{nystrup2020regime} survey machine learning approaches to regime detection, documenting improved performance over traditional econometric methods.

However, these approaches detect regimes through statistical properties of \textit{observable outcomes} (volatility clustering, return distributions) rather than the structural market mechanics that \textit{produce} those outcomes. A Markov-switching model might identify a ``high-volatility regime'' but cannot explain \textit{why} volatility increased---whether from dealer hedging, macro uncertainty, or liquidity withdrawal. Our contribution detects regimes through \textit{dealer positioning constraints}---a microstructure-based signal with explicit causal interpretation.

\subsection{Obfuscation Testing}

This article is the extended journal version of two of the authors' prior conference papers, consolidated into a single multi-scale validation framework. \citet{regan2025obfuscation} (IEEE BigData 2025, LLM-Finance Workshop) introduced obfuscation testing at the \emph{single-day} scale, demonstrating 71.5\% detection of single-day dealer constraints with 91.2\% predictive accuracy. \citet{regan2026regimes} (AIAI 2026, Springer IFIP AICT) introduced the \emph{30-day multi-day} regime-detection framework and the 2020--2025 0DTE-evolution analysis. The present article unifies these two temporal scales and adds material not present in either conference paper---a Markov-switching benchmark, a 45-configuration threshold-sensitivity analysis, bootstrap and Wilson confidence intervals throughout, and an expanded methodological appendix; the underlying experiments are those reported in the conference papers, and no new experiments were conducted for this article. \citet{ribeiro2020beyond} introduced behavioral testing for NLP models, but no prior work has developed validation methods specifically for structural reasoning in financial markets. Obfuscation testing fills this gap by removing all memorizable context while preserving mechanical relationships.

We extend this methodology to \textbf{temporal domains} (30-day regimes vs single-day snapshots), with the key challenge being selectivity: discriminating persistent regimes from transitional periods validates structural reasoning, whereas detecting universal patterns proves only pattern matching.

\subsection{Research Gap}

Despite extensive literature on dealer hedging mechanics \citep{ni2005stock,garleanu2009demand}, 0DTE market structure \citep{dim2023odtes,adams2025zerodte}, and LLM financial reasoning \citep{lopez2023can,lopezlira2025memorization}, no prior work has:
\begin{enumerate}
\item demonstrated that raw options chain data outperforms parametric GEX for structural pattern detection,
\item validated detection through obfuscation testing eliminating training data contamination,
\item established detection-alpha orthogonality proving mechanism identification independent of profitability, or
\item extended structural reasoning validation from single-day to persistent multi-day regimes.
\end{enumerate}
This work addresses all four gaps within a unified framework.
